\section{Introduction}
\label{sec:introduction}

Ordinary coends package dinatural families into a single universal object.
They support the calculus of Kan extensions, profunctors, density, and the
co-Yoneda lemma.  For a category varying over a base, however, there are two
different operations that one might call a coend.  One can form an ordinary
coend separately in each fibre, or one can ask for a universal object that is
compatible with every change of base.  The first operation is familiar but
does not, by itself, produce the second.  The elementary example in
\cref{ex:fibrewise-failure} makes this distinction explicit.

This paper begins a coend calculus adapted to categories indexed over a
Grothendieck site.  Its basic move is to formulate the universal property over
the whole slice site \(\C/U\).  A relative cowedge with apex \(X\) therefore
has a component not only for every object in \(\D(U)\), but for every
generalized object \(A\in\D(V)\) above every arrow \(v:V\to U\).  A second
axiom compares these components after further pullback.  The resulting
universal object records change-of-base data that a fibrewise definition
omits.  It becomes descent-sensitive only when the indexed inputs satisfy
the relevant \(J\)-descent hypotheses; the topology does not enter the two
defining equations by itself.

The motivation comes from Drieux's lecture series on indexed categories and
relative topos theory \cite{Drieux2025LIX,Drieux2025Wuhan,Drieux2026LHC,
Drieux2026TACL}.  Those slides give an indexed Yoneda construction, formulate
indexed colimits by localization to every slice, and announce relative
(co-)Yoneda and indexed-(co)limit projects.  They are programme-level sources,
not published proofs of the claims developed here.  We therefore reconstruct
the definitions from their universal properties and prove independently every
existence or computation statement used below.

There is important prior art.  Shulman's indexed-enriched formalism already
constructs a canceling tensor product and the composition of indexed
profunctors by a fibrewise coequalizer under indexed coproduct and
Beck--Chevalley hypotheses \cite[Definition~2.4, Definition~2.10, and
Lemma~3.25]{Shulman2013}.  Earlier, Betti and Walters developed an end calculus
for categories enriched in the bicategory of spans of a base topos
\cite{BettiWalters1989}.  Lee develops
indexed profunctors over 2-categories \cite{Lee2023}.  The present definition
is narrower in one direction and different in another: it is initially
unenriched and tied to a fixed site, but it makes localization over every
slice part of the cowedge itself.  Thus the intended new problem is not
merely to decorate the ordinary coend sign.  It is to relate coend calculus
to topology, descent, and change of site in relative topos theory.

Pitts gives an earlier and particularly close warning against claiming
indexed Fubini itself as new.  In Sections~1--2 of his 1985 paper, the words
\emph{category}, \emph{functor}, and \emph{natural transformation} mean
\(S\)-indexed ones; Lemma~1.1 uses an indexed coend construction, and the
proof of Theorem~2.3 explicitly invokes Fubini for iterated coends
\cite[Convention and Lemma~1.1, pp.~46--47; proof of Theorem~2.3,
p.~51]{Pitts1985}.  The candidate contribution here is therefore narrower:
the full-slice cowedge with its base relations, its representing
coequalizer, and a proof that this presentation is stable under the stated
pullback hypotheses.

\subsection*{Main results}

After fixing size and variance conventions, we:
\begin{enumerate}
  \item define localized relative cowedges, coends, wedges, and ends;
  \item identify a relative coend with a representing object for the
        cowedge functor;
  \item construct restriction along every map of the site and the canonical
        base-change comparison between local coends;
  \item prove invariance under a chosen indexed equivalence and reduction to
        the ordinary definition over the terminal site; and
  \item exhibit a two-object indexed category for which fibrewise coends fail
        to assemble, while the localized coends are stable;
  \item construct every fixed-base relative coend by an explicit
        coequalizer, using only the displayed coproducts, coequalizers, and
        left adjoints to reindexing;
  \item prove stable base change when the base has chosen pullbacks,
        Beck--Chevalley holds, and reindexing preserves those colimits; and
  \item prove a Fubini theorem whenever the inner relative coends exist and
        are stable under base change.
\end{enumerate}
All these statements are proved in this paper.  We do \emph{not} yet claim a
relative co-Yoneda lemma, a localization-sensitive profunctor bicategory, or
invariance under a change of site presentation.  Those are separated into
explicit problems in \cref{sec:next}.


\section{The ambient indexed data}
\label{sec:ambient}

\subsection{Sites and relative-topos conventions}
\label{subsec:sites}

Fix Grothendieck universes \(\mathcal U\in\mathcal W\) and a
\(\mathcal U\)-small site \((\C,J)\).  In this paper \(\Cat\) denotes
\(\mathcal U\)-small categories, while \(\CAT\) allows
\(\mathcal W\)-small categories with \(\mathcal U\)-small hom-sets.  Thus
the indexing collections and the sets of cowedge components below are
\(\mathcal U\)-small; the coefficient categories may have a
\(\mathcal W\)-small class of objects.  Write
\[
  \E=\Sh(\C,J),
  \qquad
  \ell_J=a_J\yon:\C\longrightarrow\E
\]
for the associated sheaf topos and the sheafified Yoneda functor.  The
canonical stack attached to this presentation is the pseudofunctor
\[
  \Sstack_{(\C,J)}:\C^{\op}\longrightarrow\CAT,
  \qquad
  U\longmapsto \E/\ell_J(U),
\]
whose reindexing functors are pullbacks.  Caramello and Zanfa prove that this
is a \(J\)-stack \cite[Definition~2.6.4 and Theorem~2.6.8]{CaramelloZanfa2021}.
It will be the principal coefficient system for applications, although the
definition below works for a general indexed category of coefficients.

In the terminology of relative topos theory, a relative topos over \(\E\) is
a geometric morphism \(p:\mathcal F\to\E\)
\cite[Definition~8.2.2]{CaramelloZanfa2021}.  We keep this distinct from an
indexed category \(\D\): the latter is input from which relative categorical
constructions are formed, not by itself a relative topos.

The smallness assumption is a working convention.  A small-generated site is
locally small and has a small \(J\)-dense subcategory
\cite[p.~4]{CaramelloZanfa2021}.  After choosing such a subcategory one may
apply the same formulas to the resulting small presentation; independence of
that choice is a separate comparison theorem and is not assumed below.  This
distinction is essential: a formula on a presentation is not automatically
an invariant of the base topos.

\begin{remark}[Where the topology enters]
\label{rem:topology-entry}
The equations in \cref{def:relative-cowedge} use the underlying category
\(\C\), not the covering families of \(J\) directly.  The topology enters
when \(\D\) and \(\V\) carry \(J\)-descent, when \(H\) respects that descent,
or through the choice \(\V=\Sstack_{(\C,J)}\).  Without such hypotheses the
construction is localization-sensitive but not yet topology-invariant.
\end{remark}

\subsection{Indexed categories, opposites, and coefficients}
\label{subsec:indexed-data}

Let
\[
  \D:\C^{\op}\longrightarrow\Cat,
  \qquad
  \V:\C^{\op}\longrightarrow\CAT
\]
be pseudofunctors.  We assume that each \(\D(U)\) is small and that each
\(\V(U)\) is locally small.  For an arrow \(u:U'\to U\), both reindexing
functors are denoted \(u^*\); the ambient pseudofunctor will always be clear.
The coherence isomorphism is written
\[
  \alpha_{v,u,X}:v^*u^*X\xrightarrow{\ \cong\ }(uv)^*X.
\]
We suppress unit constraints when no ambiguity can result.

Let \(\D^{\vee}\) denote the fibrewise opposite indexed category: its fibre
at \(U\) is \(\D(U)^{\op}\), with the induced pseudofunctorial coherence.
This is the indexed opposite used in relative category theory
\cite[Definition~2.1.4]{CaramelloZanfa2021}.  Fix an indexed functor
\[
  H:\D^{\vee}\times_{\C}\D\longrightarrow\V.
\]
Thus each fibre has a functor
\[
  H_U:\D(U)^{\op}\times\D(U)\longrightarrow\V(U),
\]
and every \(u:U'\to U\) has an invertible comparison
\begin{equation}
  \chi^H_{u,A,B}:u^*H_U(A,B)\xrightarrow{\ \cong\ }
  H_{U'}(u^*A,u^*B),
  \label{eq:H-comparison}
\end{equation}
natural in \(A,B\) and coherent under composition.  Requiring \(\chi^H\) to
be invertible is the point at which \(H\) is genuinely indexed rather than
merely laxly indexed.

For \(U\in\C\), the localization of \(\D\) at \(U\) is
\[
  \D/U:(\C/U)^{\op}\longrightarrow\Cat,
  \qquad
  (v:V\to U)\longmapsto\D(V).
\]
If \(X\in\V(U)\), its cartesian constant section over \(\C/U\) has value
\(v^*X\) at \(v:V\to U\).  These two constructions supply the domain and
apex of the relative cowedges below.


\section{Localized cowedges and relative coends}
\label{sec:definition}

Recall that for an ordinary functor \(K:\mathcal A^{\op}\times\mathcal
A\to\mathcal B\), a cowedge to \(X\) is a dinatural family
\(K(A,A)\to X\); its initial cowedge is the coend
\cite[Definition~1.1.4 and Definition~1.1.6]{Loregian2023}.  We add a base
coherence condition to this familiar definition.

\begin{definition}[Localized relative cowedge]
\label{def:relative-cowedge}
Fix \(U\in\C\) and \(X\in\V(U)\).  A \emph{relative cowedge from \(H\) to
\(X\) over \(U\)} is a family
\begin{equation}
  \omega_{v,A}:H_V(A,A)\longrightarrow v^*X
  \label{eq:cow-component}
\end{equation}
for every arrow \(v:V\to U\) and every \(A\in\D(V)\), satisfying:
\begin{enumerate}[label=(C\arabic*)]
  \item \emph{fibrewise dinaturality}: for each \(f:A\to B\) in
  \(\D(V)\),
  \begin{equation}
    \omega_{v,A}\,H_V(f^{\op},1_A)
    =
    \omega_{v,B}\,H_V(1_B,f)
    :H_V(B,A)\longrightarrow v^*X;
    \label{eq:dinaturality}
  \end{equation}
  \item \emph{base coherence}: for every \(w:W\to V\),
  \begin{equation}
    \omega_{vw,w^*A}\,\chi^H_{w,A,A}
    =
    \alpha_{w,v,X}\,w^*(\omega_{v,A})
    :w^*H_V(A,A)\longrightarrow(vw)^*X.
    \label{eq:base-coherence}
  \end{equation}
\end{enumerate}
The set of such families is denoted \(\RCow_U(H;X)\).
\end{definition}

The definition ranges over the Grothendieck construction of \(\D/U\), but
writing out \cref{eq:dinaturality,eq:base-coherence} avoids choosing a
cleavage for that construction.  The two equations have different jobs:
dinaturality moves within one fibre, whereas base coherence moves between
fibres.

A morphism \((X,\omega)\to(Y,\tau)\) is a morphism \(m:X\to Y\) in
\(\V(U)\) such that
\[
  v^*m\,\omega_{v,A}=\tau_{v,A}
\]
for all \(v,A\).  This defines the category \(\Cow_U(H)\).  Equivalently,
\(\Cow_U(H)\) is the category of elements of the covariant functor
\[
  \RCow_U(H;-):\V(U)\longrightarrow\Set.
\]

\begin{definition}[Relative coend]
\label{def:relative-coend}
A \emph{relative coend of \(H\) over \(U\)} is an initial object
\((L_U,\lambda^U)\) of \(\Cow_U(H)\).  When it exists, we write
\[
  L_U=\mathop{\int^{A\in\D}_{\Rel,U}} H(A,A).
\]
The notation records that the cowedge is localized over \(\C/U\); it is not
an assertion that an ordinary fibrewise coend has already been computed.
\end{definition}

\begin{proposition}[Representing property and uniqueness]
\label{prop:representing}
For \(L_U\in\V(U)\), the following are equivalent.
\begin{enumerate}
  \item There is a relative cowedge \(\lambda^U\) making
        \((L_U,\lambda^U)\) initial in \(\Cow_U(H)\).
  \item There is a natural isomorphism
  \begin{equation}
    \V(U)(L_U,X)\cong\RCow_U(H;X)
    \label{eq:representing}
  \end{equation}
  sending \(1_{L_U}\) to \(\lambda^U\).
\end{enumerate}
Consequently, a relative coend is unique up to a unique isomorphism
compatible with its universal cowedge.
\end{proposition}

\begin{proof}
Given an initial cowedge, send \(m:L_U\to X\) to the family
\((v^*m\,\lambda^U_{v,A})_{v,A}\).  Initiality gives the inverse and its
uniqueness, while composition in \(\V(U)\) gives naturality in \(X\).
Conversely, the image of \(1_{L_U}\) under \cref{eq:representing} is a
universal element, hence an initial cowedge.  The final assertion is the
standard uniqueness of representing objects, including compatibility with
the distinguished universal element.
\end{proof}

\begin{definition}[Relative wedge and end]
\label{def:relative-end}
For an indexed functor \(K:\D^{\vee}\times_{\C}\D\to\V\), a relative wedge
from \(X\in\V(U)\) to \(K\) is a family
\[
  \kappa_{v,A}:v^*X\longrightarrow K_V(A,A)
\]
for \(v:V\to U\) and \(A\in\D(V)\), subject to two equations.  For
\(f:A\to B\) in \(\D(V)\),
\begin{equation}
  K_V(1_A,f)\,\kappa_{v,A}
  =K_V(f^{\op},1_B)\,\kappa_{v,B}
  :v^*X\longrightarrow K_V(A,B),
  \label{eq:wedge-dinaturality}
\end{equation}
and for \(w:W\to V\),
\begin{equation}
  \kappa_{vw,w^*A}\,\alpha_{w,v,X}
  =\chi^K_{w,A,A}\,w^*\kappa_{v,A}
  :w^*v^*X\longrightarrow K_W(w^*A,w^*A).
  \label{eq:wedge-base-coherence}
\end{equation}
A morphism \((X,\kappa)\to(Y,\tau)\) is a map \(m:X\to Y\) with
\(\kappa_{v,A}=\tau_{v,A}v^*m\) for every \(v,A\).  These objects and
morphisms form \(\operatorname{Wedge}_U(K)\).  A \emph{relative end} is a
terminal object of \(\operatorname{Wedge}_U(K)\), written
\[
  \mathop{\int_{A\in\D}^{\Rel,U}}K(A,A).
\]
No completeness of the fibrewise opposite coefficient system is assumed.
\end{definition}


\section{Restriction and stable base change}
\label{sec:base-change}

The use of all generalized objects has an immediate payoff: cowedges restrict
without any existence assumption.

\begin{proposition}[Restriction to a slice]
\label{prop:restriction}
Every arrow \(u:U'\to U\) induces a functor
\[
  \rho_u:\Cow_U(H)\longrightarrow\Cow_{U'}(H).
\]
It sends an apex \(X\) to \(u^*X\).  Up to the pseudofunctorial constraints,
the component of \(\rho_u\omega\) at \(v:V\to U'\) and
\(A\in\D(V)\) is \(\omega_{uv,A}\).  These functors are pseudofunctorial:
for \(U''\xrightarrow{v}U'\xrightarrow{u}U\), the compositor
\[
  \gamma_{v,u}:\rho_v\rho_u\Longrightarrow\rho_{uv}
\]
has apex component
\(\alpha_{v,u,X}:v^*u^*X\to(uv)^*X\), and the unitor is the unit
constraint of \(\V\).
\end{proposition}

\begin{proof}
Define
\[
  (\rho_u\omega)_{v,A}:
  H_V(A,A)\xrightarrow{\omega_{uv,A}}(uv)^*X
  \xrightarrow{\alpha^{-1}_{v,u,X}}v^*u^*X.
\]
The dinaturality equation is exactly that for \(\omega_{uv,-}\).  Base
coherence follows by pasting the base-coherence square for \(\omega\) with
the associativity coherence for \(\alpha\).  On morphisms put
\(\rho_u(m)=u^*m\).  For the compositor, the defining equality for a
morphism of cowedges is the associativity pentagon for \(\alpha\), applied
to \(\omega_{uvz,A}\); hence its apex is exactly
\(\alpha_{v,u,X}\).  The unitor check is the triangle axiom.  The pentagon
and triangle for \(\gamma\) are therefore those of the pseudofunctor
\(\V\), component by component.
\end{proof}

Suppose relative coends \((L_U,\lambda^U)\) have been chosen for all
\(U\in\C\).  The restricted cowedge \(\rho_u\lambda^U\) has apex \(u^*L_U\).
Initiality at \(U'\) therefore supplies a unique comparison
\begin{equation}
  b_u:L_{U'}\longrightarrow u^*L_U
  \label{eq:base-change-map}
\end{equation}
such that \(b_u\lambda^{U'}=\rho_u\lambda^U\), with the structural
isomorphisms understood.

\begin{definition}[Stable relative coends]
\label{def:stable}
A chosen family of relative coends is \emph{stable under base change} if
every canonical map \(b_u\) in \cref{eq:base-change-map} is an isomorphism.
\end{definition}

\begin{theorem}[Coherent assembly]
\label{thm:assembly}
The canonical maps \(b_u\) satisfy the identity and composition laws up to
the coherence of \(\V\).  If the relative coends are stable, the inverses
\[
  b_u^{-1}:u^*L_U\xrightarrow{\ \cong\ }L_{U'}
\]
make \((L_U)_{U\in\C}\) a cartesian pseudonatural object of \(\V\).  This
structure is uniquely determined by the universal cowedges.
\end{theorem}

\begin{proof}
For an identity arrow, both the unit comparison and \(b_{1_U}\) are
morphisms from the initial cowedge to its unit restriction; they agree by
uniqueness.  For composable arrows \(U''\xrightarrow{v}U'\xrightarrow{u}U\),
the two morphisms from the initial cowedge over \(U''\) to the
\((uv)\)-restriction of \(\lambda^U\) are equal.  On apexes this says
\begin{equation}
  b_{uv}
  =\alpha_{v,u,L_U}\,v^*b_u\,b_v.
  \label{eq:b-composition}
\end{equation}
If all \(b_u\) are invertible, put
\(\phi_u=b_u^{-1}:u^*L_U\to L_{U'}\).
\Cref{eq:b-composition} is equivalent to
\begin{equation}
  \phi_v\,v^*\phi_u
  =\phi_{uv}\,\alpha_{v,u,L_U},
  \label{eq:phi-composition}
\end{equation}
and the identity argument gives its unit law.  Thus \(U\mapsto L_U\), with
the invertible cartesian transition maps \(\phi_u\), is a cartesian section
of the Grothendieck fibration of \(\V\), equivalently a pseudonatural
transformation from the terminal pseudofunctor to \(\V\).  Any other
compatible transitions would invert morphisms with the same universal
property and hence coincide with them.
\end{proof}

\begin{remark}
\label{rem:existence-vs-stability}
Existence and stability are logically separate.  The definition gives
\(b_u\) only after relative coends exist on both slices; it does not force
that map to be invertible.  A useful existence theorem must therefore either
prove stability or state it as an additional hypothesis.
\end{remark}


\section{First comparison theorems}
\label{sec:comparisons}

\begin{proposition}[Indexed-equivalence invariance]
\label{prop:indexed-equivalence}
Let \(F:\D\rightleftarrows\D':G\) be a chosen pseudonatural adjoint
equivalence of indexed categories, with invertible indexed unit
\(\eta:1_{\D}\Rightarrow GF\) and counit
\(\varepsilon:FG\Rightarrow1_{\D'}\).  Let
\(H':(\D')^{\vee}\times_{\C}\D'\to\V\), and suppose an indexed natural
isomorphism
\[
  \theta:H\xRightarrow{\ \cong\ }H'\,(F^{\vee}\times F)
\]
has been fixed.  Then for every \(U\) there is an equivalence
\[
  \Cow_U(H)\simeq\Cow_U(H')
\]
that is the identity on apex objects.  In particular, one side has a relative
coend if and only if the other does, and the resulting apex objects are
canonically isomorphic relative to the chosen equivalence data.
\end{proposition}

\begin{proof}
For an \(H'\)-cowedge \(\tau\) with apex \(X\), define \(T\tau\) by
\begin{equation}
  H_V(A,A)\xrightarrow{\theta_{V,A,A}}
  H'_V(F_VA,F_VA)\xrightarrow{\tau_{v,F_VA}}v^*X.
  \label{eq:equivalence-forward}
\end{equation}
Naturality and indexed coherence of \(\theta\) give (C1) and (C2).
Conversely, for an \(H\)-cowedge \(\omega\), define \(S\omega\) at
\(A'\in\D'(V)\) as
\begin{equation}
\begin{split}
  H'_V(A',A')
  &\xrightarrow{H'_V(\varepsilon_{V,A'}^{\op},
                         \varepsilon_{V,A'}^{-1})}
    H'_V(F_VG_VA',F_VG_VA')\\
  &\xrightarrow{\theta^{-1}_{V,G_VA',G_VA'}}
    H_V(G_VA',G_VA')
  \xrightarrow{\omega_{v,G_VA'}}v^*X.
\end{split}
\label{eq:equivalence-backward}
\end{equation}
Here the first arrow is well typed because the first variable is
contravariant.  Naturality of \(\varepsilon\) proves (C1), while the fact
that \(\varepsilon\) and \(\theta\) are indexed modifications proves (C2).
Both \(S\) and \(T\) leave apexes and apex morphisms unchanged.

At an object \(A\in\D(V)\), the component of \(TS\omega\) is obtained from
\(\omega_{v,G_VF_VA}\) along \(\eta_{V,A}:A\to G_VF_VA\).  Dinaturality of
\(\omega\), followed by the triangle identity
\(\varepsilon_{F A}F\eta_A=1_{F A}\), identifies it with
\(\omega_{v,A}\).  The same calculation for \(A'\in\D'(V)\), now using
\(G\varepsilon_{A'}\eta_{G A'}=1_{G A'}\), identifies \(ST\tau\) with
\(\tau\).  These componentwise identifications are indexed by coherence of
\(\eta,\varepsilon,\theta\), and are the identity on apexes.  Thus \(S\)
and \(T\) are inverse equivalences of cowedge categories.  Initial objects
correspond, and their apexes are uniquely isomorphic by
\cref{prop:representing}.
\end{proof}

The word \emph{chosen} in \cref{prop:indexed-equivalence} is deliberate.  An
unstructured statement that the fibres merely happen to be equivalent does
not provide the coherence needed to transport \cref{eq:base-coherence}.

\begin{proposition}[Terminal-site reduction]
\label{prop:terminal-reduction}
Let \(\C=\one\) with its trivial (coarsest) Grothendieck topology.  If
\(\D(*)=\mathcal A\), \(\V(*)=\mathcal B\), and
\(H_*:\mathcal A^{\op}\times\mathcal A\to\mathcal B\), then the
relative-cowedge category over \(*\) is canonically isomorphic, through the
pseudofunctor unit constraints, to the ordinary cowedge category for \(H_*\).
Consequently,
\[
  \mathop{\int^{A\in\D}_{\Rel,*}}H(A,A)
  \cong
  \int^{A\in\mathcal A}H_*(A,A)
\]
whenever either side exists.
\end{proposition}

\begin{proof}
The slice \(\one/*\) has one object and one arrow.  Thus
\cref{eq:cow-component} is an ordinary family indexed by
\(A\in\mathcal A\), \cref{eq:dinaturality} is ordinary dinaturality, and
\cref{eq:base-coherence} reduces to the unit coherence.  Normalizing the
pseudofunctors makes the two definitions literal; without normalization,
the unit constraints give the asserted canonical isomorphism.
\end{proof}

\begin{example}[Fibrewise coends do not automatically assemble]
\label{ex:fibrewise-failure}
Let \(\C\) be the category \(0\xrightarrow{u}1\), equipped with the trivial
topology, and let \(\V\) be the constant indexed category \(\Set\).  Define
\[
  \D(1)=\varnothing,
  \qquad
  \D(0)=\one,
\]
with the unique reindexing functor \(u^*:\varnothing\to\one\).  Let
\(H_1\) be the unique functor with empty domain and let
\(H_0(*,*)=\{*\}\).

The ordinary fibrewise coend is \(\varnothing\) over \(1\) and a singleton
over \(0\).  Since \(u^*\) on \(\Set\) is the identity, these two objects
cannot be the values of a cartesian indexed object: there is not even a map
\(\{*\}\to\varnothing\) in the direction of \cref{eq:base-change-map}.

The localized relative coends behave differently.  Over \(0\), a relative
cowedge to \(X\) is one map \(\{*\}\to X\), hence a pointed set.  Over \(1\),
the fibre \(\D(1)\) contributes no component, but the generalized object
\(*\in\D(0)\) above \(u:0\to1\) again contributes one point of \(X\).
Thus \(\Cow_0(H)\) and \(\Cow_1(H)\) are both the category of pointed sets.
Their initial objects are singletons, and the base-change map is the identity.
In this example the localization-sensitive definition restores stable
assembly precisely by remembering the object visible after pullback.
\end{example}


\section{Interface with relative topos theory and prior calculus}
\label{sec:interface}

\subsection{The canonical stack as coefficients}

Taking \(\V=\Sstack_{(\C,J)}\) places the apex of the relative coend over
\(U\) in the slice topos \(\E/\ell_J(U)\).  Pullback supplies the reindexing
used in \cref{eq:base-coherence}, so a chosen stable family of relative
coends determines a cartesian section of the Grothendieck fibration of the
canonical stack.  This is a direct fixed-site link to relative topos theory,
not yet an intrinsic relative-topos construction.  A further
identification of such objects with a construction
intrinsic to an arbitrary geometric morphism \(\mathcal F\to\E\) requires a
descent theorem and is not used here.

Drieux's indexed hom construction is the motivating special case.  In the
notation of the slides, its value above \(U\) has generalized components
\begin{equation}
  [v:V\to U]\longmapsto
  \D(V)(v^*Y,v^*X),
  \label{eq:drieux-hom}
\end{equation}
followed, when required, by sheafification.  The 2026 slides then formulate an
indexed colimit through universal properties on every localization and state
an indexed co-Yoneda computation
\cite[logical slides~13--20]{Drieux2026LHC}
\cite[logical slides~13--22]{Drieux2026TACL}.  Formula
\cref{eq:drieux-hom} supports the shape of
\cref{def:relative-cowedge}; the slides do not supply a published proof of a
general relative-coend theorem, so none is inferred from them.

\subsection{Relation to Shulman's indexed profunctors}

For an \(S\)-indexed symmetric monoidal category with indexed coproducts,
Shulman defines a canceling tensor product.  Under Beck--Chevalley, tensor
preservation, and fibrewise coequalizer hypotheses, the composite of indexed
profunctors is computed by a coend-shaped coequalizer
\cite[Definition~2.4, Definition~2.10, and Lemma~3.25, pp.~7--8 and
21--22]{Shulman2013}.
This is an established indexed coend calculus and must be treated as the
principal comparison theory, not as an anticipated application of the
present paper.

The two constructions have different inputs.  Shulman's formula begins with
an indexed monoidal base and packages fibrewise coproducts together with
change-of-base adjoints.  Definition~\ref{def:relative-cowedge} instead begins
with an indexed functor \(H\) and enlarges each cowedge over \(U\) to every
object of \(\C/U\).  A substantive comparison theorem should identify
hypotheses under which Shulman's canceling tensor represents the localized
cowedge functor.  Until that theorem is proved, the two constructions are
related models rather than interchangeable notation.

Betti--Walters is an equally important novelty boundary: the title and
published abstract of their paper identify an end calculus over a base topos
in the span-enriched setting \cite{BettiWalters1989}.  The construction here
is instead written on a site presentation and builds generalized objects in
every slice into the cowedge.  Since the present source audit did not recover
a full-text comparison at theorem level, no equivalence between the two
formalisms is asserted without a separate comparison theorem.

\subsection{Presentation boundary}

An internal category in \(\E\), a stack on \((\C,J)\), and a pseudofunctor
on a chosen site are connected by comparison results, but they are not
definitionally identical.  Likewise, equivalent sites can present the same
topos without giving literally equal indexing categories.  The definition in
\cref{def:relative-cowedge} is therefore a fixed-presentation construction.
Its invariance under indexed equivalence
(\cref{prop:indexed-equivalence}) is a first step, but it is not yet a Morita
or change-of-site invariance theorem.
